\documentclass[../../main.tex]{subfiles}
\graphicspath{{\subfix{./image/}}{\subfix{../../image/}}}
\begin{document}
% SOURCE p=161; printed=146; object=ch04-notes-title; region=bottom
\section{注记与参考文献}\label{section:ch04-notes-references}
% SOURCE p=161,162; printed=146,147; object=ch04-notes-topology-history; region=bottom-top
拓扑空间概念的萌芽已经明显见于 Riemann 在 1854 年所作的关于几何基础的演讲\cite{folland-bib-113}，但又过了半个世纪，数学界才准备好系统地考虑抽象空间。第一次为极限和连续性的研究建立抽象框架的尝试，是 Fréchet 在 1906 年作出的\cite{folland-bib-052}；他引入了度量空间，以及一类更一般的拟拓扑空间，其性质用序列收敛来定义。几年后，Hausdorff\cite{folland-bib-068}提出了关于点的邻域的公理，其实质就是 Hausdorff 空间的定义，并由此推导了许多一般拓扑的基本结果。他的观点的效用很快得到认识，并成为这门学科进一步发展的基础。
% SOURCE p=162; printed=147; object=ch04-notes-topology-books; region=upper
读者可参阅若干优秀著作，获得更全面的点集拓扑论述，包括 Bourbaki\cite{folland-bib-020}、Dugundji\cite{folland-bib-034}、Engelking\cite{folland-bib-038}、Kelley\cite{folland-bib-083}和 Nagata\cite{folland-bib-102}。Engelking\cite{folland-bib-038}包含广泛的参考文献与历史注记。
% SOURCE p=162; printed=147; object=ch04-notes-separation-history; region=upper-middle
\cref{section:ch04-continuous-maps}：Urysohn 引理和 Tietze 延拓定理都最先在 Urysohn\cite{folland-bib-152}中得到证明。此前，包括 Tietze 在内的若干作者已经得到后一定理的特殊情形，相关文献见\cite{folland-bib-152}。完全正则但不正规的空间，以及正则但不完全正则的空间，其例子都相当复杂，最先由 Tychonoff\cite{folland-bib-151}构造。特别值得注意的是：存在一个正则空间，其上没有非常值连续函数；这一结果属于 Hewitt\cite{folland-bib-073}。上面列出的著作中也能找到例子。
% PROOF-001 construction is authored below; independent rounds and final verification remain required.
\Needspace{6\baselineskip}
\subfile{supplement-regular-constant.tex}
% SOURCE p=162; printed=147; object=ch04-notes-net-filter-history; region=middle
\cref{section:ch04-nets}：网的理论有时称为 Moore--Smith 收敛理论，以其创始者命名\cite{folland-bib-101}。另一种一般收敛理论由 H. Cartan 创立，经 Bourbaki 传播，以滤子的概念为基础。
% SOURCE p=162; printed=147; object=ch04-filter-definition; region=middle
\begin{definition}[滤子]\label{definition:ch04-filter}
集合 $X$ 中的一个\textbf{滤子}是具有下列性质的族 $\mathcal F\subseteq\mathcal P(X)$：
\begin{enumerate}
\item\label{item:ch04-filter-upward} 若 $F\in\mathcal F$ 且 $F\subseteq E\subseteq X$，则 $E\in\mathcal F$。
\item\label{item:ch04-filter-intersection} 若 $E,F\in\mathcal F$，则 $E\cap F\in\mathcal F$。
\item\label{item:ch04-filter-proper} $\varnothing\notin\mathcal F$。
\end{enumerate}
\end{definition}
% SOURCE p=162; printed=147; object=ch04-filter-convergence-definition; region=middle-lower
\begin{definition}[滤子的收敛]\label{definition:ch04-filter-convergence}
若 $X$ 为拓扑空间，$X$ 中的滤子 $\mathcal F$ \textbf{收敛}到 $x\in X$，是指 $x$ 的每个邻域都属于 $\mathcal F$。
\end{definition}
% SOURCE p=162; printed=147; object=ch04-derived-filter-associated-net; region=lower
\begin{definition}[导出滤子与关联的网]\label{definition:ch04-derived-filter-associated-net}
若 $\{x_\alpha\}_{\alpha\in A}$ 是 $X$ 中的网，则它的\textbf{导出滤子}是所有满足“$\{x_\alpha\}$ 最终进入 $E$”的集合 $E\subseteq X$ 所组成的族。

反过来，若 $\mathcal F$ 为滤子，则它按反向包含关系成为有向集。以 $\mathcal F$ 为指标的网 $\{x_F\}_{F\in\mathcal F}$，若对所有 $F\in\mathcal F$ 都有 $x_F\in F$，则称它与 $\mathcal F$ \textbf{关联}。
\end{definition}
% SOURCE p=162; printed=147; object=ch04-net-filter-equivalence; region=lower
\begin{proposition}\label{proposition:ch04-net-filter-equivalence}
网 $\{x_\alpha\}$ 收敛到 $x$，当且仅当其导出滤子收敛到 $x$；滤子 $\mathcal F$ 收敛到 $x$，当且仅当与它关联的所有网都收敛到 $x$。
\end{proposition}
% SOURCE p=162; printed=147; object=ch04-net-filter-proof-reference; region=lower
这些事实很容易验证。更多内容见 Bourbaki\cite{folland-bib-020}或 Dugundji\cite{folland-bib-034}。
% SUPPLEMENT: id=v14-ch04-filter-net-details; parent=ch04-net-filter-equivalence; kind=omitted-proof+source-correction+motivation+hypotheses; origin=independent-derivation
\begin{remark}
\textbf{非空约定与定义的动机（重构新增）。}原列出的三条二元条件未明说 $\mathcal F\ne\varnothing$；按字面，空族也满足三条。通常的滤子定义还要求族非空，等价地在向上封闭条件下要求 $X\in\mathcal F$。后续把它当作非空有向集，必须补入这一约定。以下均如此理解；不改动上面的原三条。

滤子记录“最终应当进入哪些集合”：若最终进入一个集合，也最终进入其超集；若分别最终进入两个集合，也最终进入其交；不可能最终进入空集。邻域全部属于滤子便表示它具有趋近指定点的全部局部要求。反向包含把更小、约束更强的集合放在更后的指标位置。关联的网选择每个约束集合中的一个代表点，选择可以有很多种，因此后面的判别要求\textbf{所有}关联的网收敛。
\end{remark}
\begin{proof}
\textbf{良定义与等价性的完整补写（重构新增）。}先设网的指标集 $A$ 非空且有向，记导出族为 $\mathcal D$。整个 $X\in\mathcal D$，所以 $\mathcal D$ 非空；向上封闭直接由“最终进入”的定义得到。如果网从 $\alpha_0$ 起在 $E$ 中、从 $\alpha_1$ 起在 $F$ 中，取共同上界 $\alpha_2\ge\alpha_0,\alpha_1$，则从 $\alpha_2$ 起在 $E\cap F$ 中。由于指标尾部至少含其起始指标本身，网不可能最终进入空集。故 $\mathcal D$ 确为滤子。

再设 $\mathcal F$ 为非空滤子，定义 $F\preceq G$ 当且仅当 $G\subseteq F$。这关系自反、传递；任给 $E,F\in\mathcal F$，交 $E\cap F\in\mathcal F$ 是二者在该关系下的共同上界，所以为有向集。每个成员都非空，选择公理可给出每个 $F$ 中的代表 $x_F$，故关联的网存在。

网收敛到 $x$，按定义等价于它最终进入 $x$ 的每个邻域；这又恰好等价于每个邻域属于 $\mathcal D$，即 $\mathcal D\to x$。

若 $\mathcal F\to x$，任取关联的网 $(x_F)$ 及 $x$ 的邻域 $U$。因为 $U\in\mathcal F$，当 $F\succeq U$ 时 $F\subseteq U$，所以 $x_F\in U$。因此该网最终进入每个邻域，即收敛到 $x$。

反过来，若 $\mathcal F$ 不收敛到 $x$，有 $x$ 的邻域 $U\notin\mathcal F$。对每个 $F\in\mathcal F$，必有 $F\setminus U\ne\varnothing$，否则 $F\subseteq U$ 及向上封闭将给 $U\in\mathcal F$。从每个非空 $F\setminus U$ 选取 $x_F$，得到一个关联的网，它从不进入 $U$，所以不收敛到 $x$。这否定“所有关联的网都收敛”，证明最后的反向。
\end{proof}
\begin{remark}
\textbf{逐条件说明（重构新增）。}有限交封闭对应网的有向性，只需给有限多个最终要求找共同起点；不要求无限交仍属于滤子，否则通常的邻域和尾部滤子会被排除。排除空集确保每个约束都可满足，族非空确保有指标可以开始。证明关联的网存在及反向构造不收敛的网时，明确使用了选择公理；这与本书的通常集合论背景一致。空空间不存在这种非空且不含空集的滤子，也不存在取值于其中的非空指标网。
\end{remark}
% END SUPPLEMENT: id=v14-ch04-filter-net-details
% SOURCE p=162; printed=147; object=ch04-notes-compact-terminology; region=bottom
\cref{section:ch04-compact-spaces}：“紧”的用法并不完全统一。在许多较早的著作中，术语 compact 和 bicompact 分别指可数紧和紧；有些作者则用 compact 和 quasi-compact 分别指紧 Hausdorff 和紧。文献中常用 conditionally compact（条件紧）和 relatively compact（相对紧）作为 precompact（预紧）的同义词；后一个名称不太恰当，因为容易让人以为它指相对拓扑下的紧性，而二者相差很大。
% SUPPLEMENT: id=v14-ch04-compact-terminology-remark; parent=ch04-notes-compact-terminology; kind=motivation+counterexample; origin=independent-derivation
\begin{remark}
\textbf{术语区别的例证（重构新增）。}在 $\mathbb R$ 中，$(0,1)$ 的闭包 $[0,1]$ 紧，所以它按本书定义预紧；但 $(0,1)$ 取相对拓扑并不紧。例如 $\{(1/n,1):n\ge2\}$ 是其相对开覆盖，任意有限子族的并仍漏掉足够接近 $0$ 的点。因此“闭包在环境空间中紧”与“子集自身在相对拓扑中紧”不能混用。
\end{remark}
% END SUPPLEMENT: id=v14-ch04-compact-terminology-remark
% SOURCE p=163; printed=148; object=ch04-notes-tychonoff-history; region=top
\cref{section:ch04-tychonoff}：Tychonoff\cite{folland-bib-151}证明了对任意集合 $A$，$[0,1]^A$ 都紧；结合也在同一篇论文中的\cref{corollary:ch04-compact-closed-cube}，很容易推出任意族紧 Hausdorff 空间的乘积紧。完全一般的 Tychonoff 定理属于 Čech\cite{folland-bib-023}。我们给出的证明比早期证明更简洁，出自 Chernoff\cite{folland-bib-024}。
% SUPPLEMENT: id=v14-ch04-cube-product-compact-proof; parent=ch04-notes-tychonoff-history; kind=omitted-proof+motivation+hypotheses; origin=independent-derivation
\begin{proof}
\textbf{从立方体到紧 Hausdorff 乘积的补写（重构新增）。}设各 $X_\alpha$ 紧 Hausdorff。若某个因子为空，积为空而紧；空指标积为单点，也紧。其余情形对每个 $\alpha$ 用全体 $I$ 值连续函数 $\mathcal F_\alpha=C(X_\alpha,I)$ 作典范嵌入，得闭像 $Z_\alpha\subseteq I^{\mathcal F_\alpha}$。令 $B$ 为所有标记坐标 $(\alpha,f)$（$f\in\mathcal F_\alpha$）的集合。

把各组坐标合在一起，$\prod_\alpha I^{\mathcal F_\alpha}$ 与 $I^B$ 同胚：两个方向都只重新排列坐标，因而逐坐标连续。$\prod_\alpha Z_\alpha$ 对应于
\begin{align*}
\bigcap_\alpha r_\alpha^{-1}(Z_\alpha),
\end{align*}
其中 $r_\alpha:I^B\to I^{\mathcal F_\alpha}$ 是保留第 $\alpha$ 组坐标的连续映射。各 $Z_\alpha$ 闭，所以该交是紧立方体 $I^B$ 的闭子集，因而紧。各嵌入的乘积给 $\prod_\alpha X_\alpha$ 到此闭子集的同胚，结论得证。
\end{proof}
\begin{remark}
\textbf{条件说明（重构新增）。}这里从“所有立方体紧”出发，故不是再次用一般乘积紧性证明自身。Hausdorff 条件使每个紧因子可实现为立方体的闭子集；非 Hausdorff 紧空间不能由此论证直接涵盖，这正是原文区分两个历史版本的原因。
\end{remark}
% END SUPPLEMENT: id=v14-ch04-cube-product-compact-proof
% SOURCE p=163; printed=148; object=ch04-tychonoff-implies-choice; region=upper
\begin{proposition}\label{proposition:ch04-tychonoff-implies-choice}
选择公理通常以 Zorn 引理的形式，是所有 Tychonoff 定理证明中的一个基本成分。一个有趣的事实是：反过来，Tychonoff 定理蕴含选择公理；这是 Kelley\cite{folland-bib-082}发现的。
\end{proposition}
% SOURCE p=163; printed=148; object=ch04-proof-tychonoff-implies-choice; region=upper-middle
\begin{proof}
设 $\{X_\alpha\}_{\alpha\in A}$ 是非空的一族非空集合。选一个不属于任何 $X_\alpha$ 的点 $\omega$，置 $X_\alpha^*=X_\alpha\cup\{\omega\}$，并规定其开集为 $\varnothing$、$X_\alpha$、$\{\omega\}$ 和 $X_\alpha^*$。显然 $X_\alpha^*$ 紧，因此 Tychonoff 定理推出 $X^*=\prod_{\alpha\in A}X_\alpha^*$ 紧。令 $F_\alpha=\pi_\alpha^{-1}(X_\alpha)$。

这些 $F_\alpha$ 闭，并由有限族集合的选择公理——它可以从集合论其余标准公理证明——具有有限交性质。具体地，给定有限集 $B\subseteq A$，为 $\beta\in B$ 选择 $x_\beta\in X_\beta$；则 $\bigcap_{\beta\in B}F_\beta$ 包含点 $x\in X^*$，其中 $\pi_\beta(x)=x_\beta$ 对 $\beta\in B$ 成立，$\pi_\alpha(x)=\omega$ 对 $\alpha\notin B$ 成立。由\cref{proposition:ch04-compact-fip}，$\bigcap_{\alpha\in A}F_\alpha$ 非空，而它恰好是 $\prod_{\alpha\in A}X_\alpha$。
\end{proof}
% SUPPLEMENT: id=v14-ch04-tychonoff-choice-details; parent=ch04-tychonoff-implies-choice; kind=proof-detail+motivation+hypotheses; origin=independent-derivation
\begin{proof}
\textbf{避免循环使用选择的细节（重构新增）。}对任意集合 $S=\bigcup_\alpha X_\alpha$，在通常含正则公理的集合论中 $S\notin S$，因而可取一个在 $S$ 外的集合为 $\omega$，无需作逐指标选择。$X_\alpha^*$ 的四个指定开集对任意并和有限交封闭，确为拓扑。覆盖 $\omega$ 的开集为 $\{\omega\}$ 或整个空间；覆盖非空块 $X_\alpha$ 中任一点的开集为 $X_\alpha$ 或整个空间。因此每个开覆盖至多取两个成员即可覆盖，所以紧。

即使尚未知道 $\prod_\alpha X_\alpha$ 非空，$X^*$ 已有明确的点 $\alpha\mapsto\omega$。$X_\alpha$ 在 $X_\alpha^*$ 中闭，因为其补集 $\{\omega\}$ 开；连续坐标投影使 $F_\alpha$ 闭。有限个非空集合可用有限次存在量词选择取得一个元组，或对集合族的有限大小归纳证明，不依赖一般选择公理。于是原证的每个有限交都非空，紧性的闭集有限交性质给全部交非空。

全交中的一个点是函数 $x$，对每个 $\alpha$ 满足 $x(\alpha)\in X_\alpha$，这正是所求选择函数。空指标族本来就有空选择函数，所以原证假设指标族非空不会削弱结论。
\end{proof}
\begin{remark}
\textbf{动机与条件说明（重构新增）。}添加共同的备用点，让大乘积在没有选择公理时也明确非空；闭集条件再要求每个坐标最终放弃备用点。每次只施加有限多个条件容易满足，紧性把这些有限可满足性合成全部可满足性。这里使用的是包括非 Hausdorff 紧空间的一般 Tychonoff 定理，不能把假设悄悄缩成仅紧 Hausdorff 乘积而沿用此证明。
\end{remark}
% END SUPPLEMENT: id=v14-ch04-tychonoff-choice-details
% SOURCE p=163; printed=148; object=ch04-compact-powers-imply-choice; region=middle
\begin{proposition}\label{proposition:ch04-compact-powers-imply-choice}
推广上面的论证，可以从 Tychonoff 定理的以下特殊情形推出选择公理：只要 $X$ 紧，对任意集合 $A$，$X^A$ 就紧。见 Ward\cite{folland-bib-156}。
\end{proposition}
% SUPPLEMENT: id=v14-ch04-compact-powers-choice-proof; parent=ch04-compact-powers-imply-choice; kind=omitted-proof+motivation+hypotheses; origin=independent-derivation
\begin{proof}
\textbf{独立补写证明（重构新增）。}假设每个紧空间的任意幂均紧。给定非空族非空集合 $(X_\alpha)_{\alpha\in A}$，把它们改为带标签的不交非空块 $E_\alpha=\{\alpha\}\times X_\alpha$，并添一个不属于这些块的点 $\omega$。令
\begin{align*}
Z=\{\omega\}\cup\bigcup_{\alpha\in A}E_\alpha,
\qquad
U_B=Z\setminus\bigcup_{\alpha\in B}E_\alpha
\quad(B\subseteq A\text{ 有限}).
\end{align*}
规定开集为 $\varnothing$ 及全部 $U_B$。它们确实成拓扑：$U_\varnothing=Z$，有限交满足 $U_B\cap U_C=U_{B\cup C}$。对任意非空族有限指标集 $B_j$，有 $\bigcup_jU_{B_j}=U_{\cap_jB_j}$；交集包含于其中任一有限 $B_j$，故仍有限。空族开集的并为空。

$Z$ 紧。任给其开覆盖，选一个覆盖 $\omega$ 的成员 $U_B$；它只漏掉有限多个块 $E_\alpha$（$\alpha\in B$）。对每个这样的块，取其中一点及一个覆盖该点的成员。任何非空开集都是若干整块加 $\omega$ 的并，故一旦包含块中一点便包含整个块。于是这个初始成员加有限多个后选成员覆盖 $Z$；整个过程只使用有限选择。

每个 $E_\alpha$ 在 $Z$ 中闭，因为 $Z\setminus E_\alpha=U_{\{\alpha\}}$ 开。按假设，$Z^A$ 紧，并含显式的常值点 $\omega$。其闭集
\begin{align*}
F_\alpha=\{z\in Z^A:z(\alpha)\in E_\alpha\}
\end{align*}
具有有限交性质：对有限多个坐标，在各自块中有限次选择一个元素，其余坐标全部取 $\omega$ 即可。紧性给出 $z\in\bigcap_\alpha F_\alpha$。写 $z(\alpha)=(\alpha,x_\alpha)$，则 $x_\alpha\in X_\alpha$，得到原族的选择函数。空指标族直接用空函数处理。
\end{proof}
\begin{remark}
\textbf{动机与逐条件说明（重构新增）。}这次将全部集合装进同一个紧空间的不同闭块，因而只需“同一个空间的幂紧”，不必假设不同因子的一般乘积紧。所造拓扑让有限个块之外成为一个大开集，所以紧性可仅用有限选择证明；各闭块则保留每个原集合的选择任务。此空间一般不是 Hausdorff，这与命题的完整紧空间范围一致。此补证为独立推导，未冒称复现所引论文的具体证明。
\end{remark}
% END SUPPLEMENT: id=v14-ch04-compact-powers-choice-proof
% SOURCE p=163; printed=148; object=ch04-notes-ascoli-history; region=middle-lower
Arzelà 和 Ascoli 的原始结果涉及 $\mathbb R$ 上的函数，见 Arzelà\cite{folland-bib-006}。在关于 $X,Y$ 的各种假设下，研究 $C(X,Y)$ 子集紧性的其他 Arzelà--Ascoli 定理版本，可见前述著作及 Royden\cite{folland-bib-121}。
% SOURCE p=163; printed=148; object=ch04-notes-sw-history; region=lower
\cref{section:ch04-stone-weierstrass}：Stone--Weierstrass 定理最初出现在 Stone\cite{folland-bib-144}一篇冗长而困难的论文中间。后来 Stone\cite{folland-bib-145}写了一篇大为简化的阐述，介绍该定理及一些应用，至今仍值得阅读。
% SOURCE p=163; printed=148; object=ch04-notes-cube-history; region=lower
\cref{section:ch04-cube-embeddings}：这部分材料的历史始于 Urysohn\cite{folland-bib-153}，其中用的基本就是我们概述的方法来证明度量化定理。把空间嵌入立方体的技巧隐含在这篇论文中，而 Tychonoff\cite{folland-bib-151}首先明确发展了它。Stone--Čech 紧化又隐含在后者的论文中，但最先由 Stone\cite{folland-bib-144}和 Čech\cite{folland-bib-023}明确描述并研究。
% SOURCE p=163; printed=148; object=ch04-second-countable-regular-normal; region=lower
\begin{proposition}\label{proposition:ch04-second-countable-regular-normal}
每个第二可数的正则空间都正规。
\end{proposition}
% SOURCE p=163; printed=148; object=ch04-second-countable-regular-normal-reference; region=lower
这一点不难证明，见 Kelley\cite[Lemma 4.1]{folland-bib-082}。因此，Urysohn 度量化定理中的正规性假设可以换成正则性。对于任意拓扑空间可度量化的必要充分条件也是已知的，但不像 Urysohn 定理中的条件那样容易检验；见前述著作。
% AUDIT OPEN: original reference [82, Lemma 4.1] retained; bibliography identifies [82] as Kelley's two-page choice paper. Exact external reference requires verification.
% SUPPLEMENT: id=v14-ch04-second-countable-regular-normal-proof; parent=ch04-second-countable-regular-normal; kind=omitted-proof+motivation+hypotheses; origin=independent-derivation
\begin{proof}
\textbf{补写完整证明（重构新增）。}设 $\mathcal B$ 是 $X$ 的可数基，$A,B$ 是不交闭集。若其中一个为空，用空集与 $X$ 分别分离即可。以下设二者非空。

对每个 $a\in A$，正则性给开集 $W$，使 $a\in W\subseteq\overline W\subseteq X\setminus B$；再取基元 $U$ 满足 $a\in U\subseteq W$，便有 $\overline U\cap B=\varnothing$。因此所有满足该闭包条件的基元组成可数族 $(U_n)$，覆盖 $A$。同理，所有满足 $\overline V\cap A=\varnothing$ 的基元可列作 $(V_n)$，覆盖 $B$。有限族可在末尾补空集，使两个指标统一为 $n\ge1$。

定义开集
\begin{align*}
U&=\bigcup_{n=1}^{\infty}\left(U_n\setminus\bigcup_{j=1}^n\overline{V_j}\right),\\
V&=\bigcup_{n=1}^{\infty}\left(V_n\setminus\bigcup_{j=1}^n\overline{U_j}\right).
\end{align*}
每个括号是开集减去有限个闭集，故开。各 $\overline{V_j}$ 都不遇 $A$，所以 $A\subseteq U$；同理 $B\subseteq V$。若 $x\in U\cap V$，则存在 $n,m$ 使 $x\in U_n$ 且 $x\notin\overline{V_j}$ 对 $j\le n$，同时 $x\in V_m$ 且 $x\notin\overline{U_j}$ 对 $j\le m$。若 $m\le n$，与 $x\in V_m\subseteq\overline{V_m}$ 矛盾；若 $n\le m$，与 $x\in U_n\subseteq\overline{U_n}$ 矛盾。故 $U,V$ 不交，分离了任意不交闭集。

正则空间按本书约定已经 $T_1$，所以这个闭集分离性质证明 $X$ 正规。结合\cref{theorem:ch04-urysohn-metrization}，立即得到第二可数正则空间可度量化。
\end{proof}
\begin{remark}
\textbf{动机与逐条件说明（重构新增）。}正则性先为两个闭集各造一列不碰另一侧的局部开集；它们相互之间仍可能交叠。可数性使这些开集能依次排列，每一步只减去对方前面有限个闭包，既保持开性，又使任意一对相遇的候选集由较晚的一方排除。若直接减去对方全部闭包的无限并，差集未必开，这正是有限阶段安排的作用。
\end{remark}
% END SUPPLEMENT: id=v14-ch04-second-countable-regular-normal-proof
% SOURCE p=163; printed=148; object=ch04-weaker-compactification-convention; region=bottom
\begin{definition}[“紧化”的另一种用法]\label{definition:ch04-weaker-compactification}
有时，“紧化”这个术语用来指从拓扑空间 $X$ 到某个紧空间 $Y$ 的稠密子集的连续单射 $\phi:X\to Y$，不要求 $\phi$ 为嵌入。
\end{definition}
% SOURCE p=163,164; printed=148,149; object=ch04-bohr-compactification-example; region=bottom-top
这样的“紧化”来自 $C(X)$ 中分离点、但在\cref{exercise:ch04-compactification-algebras}意义下不完全正则的子代数。一个例子是 $\mathbb R$ 上“一致殆周期”函数的代数，即由函数
\begin{align*}
f_\lambda(x)=e^{i\lambda x},\qquad\lambda\in\mathbb R,
\end{align*}
生成的代数；$\mathbb R$ 的相应“紧化”称为 Bohr 紧化。见 Folland\cite[\S4.7]{folland-bib-047}。
% SUPPLEMENT: id=v14-ch04-bohr-compactification-details; parent=ch04-bohr-compactification-example; kind=omitted-proof+source-clarification+motivation+hypotheses; origin=independent-derivation
\begin{remark}
\textbf{定义与闭包约定（重构新增）。}连续单射只保证不同点不被合并，不保证像上的相对拓扑能恢复原拓扑，因此这种用法比本节的紧化定义弱。为符合“一致殆周期”及闭函数代数的通常含义，原段的“生成”须理解为在一致范数下\textbf{闭生成}，而不只是有限代数组合。以下给出具体实现、闭包区别与非嵌入的完整验证。
\end{remark}
\begin{definition*}[本例的一致殆周期代数]\phantomsection\label{definition:ch04-uniform-almost-periodic}
\textbf{补充明确定义（重构新增）。}记 $\mathcal T$ 为所有有限和 $\sum_{j=1}^m c_j e^{i\lambda_jx}$ 组成的复代数，$c_j\in\mathbb C$、$\lambda_j\in\mathbb R$。记 $AP(\mathbb R)=\overline{\mathcal T}^{\|\cdot\|_u}\subseteq BC(\mathbb R)$；其元素称为一致殆周期函数。
\end{definition*}
\begin{proof}
\textbf{Bohr 例子的补写（重构新增）。}令 $\mathbb T=\{z:|z|=1\}$，定义
\begin{align*}
b:\mathbb R\longrightarrow\mathbb T^{\mathbb R},\qquad
\pi_\lambda(b(x))=e^{i\lambda x},\qquad
Y=\overline{b(\mathbb R)}.
\end{align*}
每个坐标连续，所以 $b$ 连续。若 $x\ne y$，取 $\lambda=\pi/(x-y)$，则 $e^{i\lambda x}/e^{i\lambda y}=-1$，故 $b(x)\ne b(y)$，所以单射。$\mathbb T$ 紧 Hausdorff，其乘积及闭子集 $Y$ 也紧 Hausdorff；$b(\mathbb R)$ 按定义稠密。这已给出较弱意义的紧化。

在 $Y$ 上，坐标函数与常数生成的复代数对共轭封闭。因为在稠密像上 $\pi_{-\lambda}=\overline{\pi_\lambda}$、$\pi_0=1$，这些恒等式由连续性延伸到 $Y$。坐标分离 $Y$ 的点，故复 Stone--Weierstrass 定理给它在 $C(Y)$ 中稠密。限制到稠密像的映射 $h\mapsto h\circ b$ 等距，且把坐标生成代数恰送到 $\mathcal T$；因此其整个像恰为 $AP(\mathbb R)$。这解释了关联代数。

证明 $b$ 不是嵌入，只需证明 $b(0)$ 的每个相对邻域都包含某个 $b(t)$，其中 $t>1$。这样的邻域含有限坐标限制
\begin{align*}
|e^{i\lambda_jt}-1|<\varepsilon
\qquad(1\le j\le m)
\end{align*}
所定义的像中基本邻域，某个 $\varepsilon>0$ 对全部有限坐标均适用。若无坐标限制，任取 $t>1$。否则取整数 $M>1$、正整数 $Q$ 使 $2\pi/Q<\varepsilon$，并考虑 $Q^m+1$ 个向量
\begin{align*}
\left(\left\{\frac{kM\lambda_1}{2\pi}\right\},\ldots,
\left\{\frac{kM\lambda_m}{2\pi}\right\}\right),
\qquad 0\le k\le Q^m,
\end{align*}
其中 $\{s\}=s-\lfloor s\rfloor\in[0,1)$。将 $[0,1)^m$ 分成边长 $1/Q$ 的 $Q^m$ 个半开小立方体。抽屉原理给两个下标 $k>l$ 的向量在同一小立方体中，所以对每个 $j$，$(k-l)M\lambda_j/(2\pi)$ 与某个整数相差小于 $1/Q$。由 $|e^{iu}-1|\le|u|$，令 $t=(k-l)M\ge M>1$，便得到全部所需不等式。

因此 $b(0)$ 的任意相对邻域都不是 $b((-1,1))$ 的子集，$b((-1,1))$ 不在 $b(\mathbb R)$ 中开。若 $b$ 为嵌入，这个原空间开集的像必须相对开，矛盾。

还可直接核对 $AP(\mathbb R)$ 不满足完全正则代数的分离要求。令 $E=\mathbb R\setminus(-1,1)$。任取 $f\in AP(\mathbb R)$ 及 $\eta>0$，选 $P=\sum_{j=1}^m c_je^{i\lambda_jx}$ 使 $\|f-P\|_u<\eta/3$。若 $C=\sum_j|c_j|>0$，用上面的有限频率构造取得 $t>1$，使各 $|e^{i\lambda_jt}-1|<\eta/(3C)$；于是
\begin{align*}
|f(t)-f(0)|
\le |f(t)-P(t)|+|P(t)-P(0)|+|P(0)-f(0)|<\eta.
\end{align*}
若 $C=0$，$P=0$，任取 $t>1$ 也给差小于 $2\eta/3$。故 $f(0)\in\overline{f(E)}$ 对所有 $f\in AP(\mathbb R)$ 成立，尤其对其中 $I$ 值函数成立，无法分离点 $0$ 与闭集 $E$。

\Needspace{20\baselineskip}
最后验证闭生成确与只作有限组合不同。级数 $h(x)=\sum_{n=1}^{\infty}2^{-n}e^{inx}$ 因一致尾和估计而属于 $AP(\mathbb R)$。对任意实数 $a$，直接积分得
\begin{align*}
\lim_{T\to\infty}\frac1{2T}\int_{-T}^T e^{iax}\,dx
=\begin{cases}1,&a=0,\\0,&a\ne0.\end{cases}
\end{align*}
对有限和可逐项使用此式。平均积分的模不超过函数的一致范数，所以对上述一致收敛级数，可先截断使尾部对所有 $T$ 的平均误差任意小，再让 $T\to\infty$。于是对每个正整数 $k$，
\begin{align*}
\lim_{T\to\infty}\frac1{2T}\int_{-T}^T h(x)e^{-ikx}\,dx=2^{-k}.
\end{align*}
然而任意有限指数和只有有限多个频率，对除这些频率外的正整数 $k$，对应平均必为零。$h$ 有无限多个非零平均系数，所以不可能为有限指数和，即 $h\notin\mathcal T$。因此此处一致闭包确不可省略。
\end{proof}
\begin{remark}
\textbf{动机与逐条件说明（重构新增）。}字符 $e^{i\lambda x}$ 能区别任意两个固定实数，但对有限多个频率，任意远处仍能同时几乎重复原点的读数。乘积邻域只检验有限坐标，所以新相对拓扑比原来的实线拓扑弱。这说明分点与分离点闭集为何是不同条件。一致闭包保持有界性、连续性、常数及共轭封闭，并使函数代数与紧目标上的全部连续函数相匹配；逐点闭包不具有本段使用的统一误差控制。
\end{remark}
% END SUPPLEMENT: id=v14-ch04-bohr-compactification-details
\end{document}
